\Askhsh{1820}{4}
\begin{ylh}{1}
    \item 5.2. Παραλληλόγραμμα
    \item 5.6. Εφαρμογές στα τρίγωνα
    \item 5.7. Βαρύκεντρο τριγώνου.
\end{ylh}
Δίνεται τρίγωνο $\mathrm{A}\mathrm{B}\Gamma$ και οι διάμεσοί του $\mathrm{A}\Delta, \mathrm{B}\mathrm{E}$ και $\Gamma\mathrm{Z}$. Προεκτείνουμε το τμήμα $\mathrm{Z}\mathrm{E}$ (προς το $\mathrm{E}$) κατά τμήμα $\mathrm{E}\mathrm{H} = \mathrm{Z}\mathrm{E}$.
Να αποδείξετε ότι:
\begin{alist}
    \item Το τετράπλευρο $\mathrm{E}\mathrm{H}\Delta\mathrm{B}$ είναι παραλληλόγραμμο. \monades{8}
    \item Η περίμετρος του τριγώνου $\mathrm{A}\Delta\mathrm{H}$ είναι ίση με το άθροισμα των διαμέσων του τριγώνου $\mathrm{A}\mathrm{B}\Gamma$. \monades{9}
    \item Οι ευθείες $\mathrm{B}\mathrm{E}$ και $\Delta\mathrm{H}$ τριχοτομούν το τμήμα $\mathrm{Z}\Gamma$. \monades{8}
\end{alist}
